\chapter{Document Ingestion and Understanding}
\label{chap:document-ingestion}

The question-answering capabilities of the application depend on a reliable
document representation. A language model cannot recover information that was
never extracted, and a retrieval algorithm cannot provide a trustworthy page
reference if source location was discarded during preprocessing. For this
reason, ingestion is not treated as a preliminary file-conversion utility. It
is an integrated pipeline that converts an authenticated upload into
searchable, traceable knowledge while preserving enough diagnostic state to
explain how that result was obtained.

The implemented path proceeds from validation through extraction, technical
PDF analysis, selective optical character recognition (OCR), normalization,
document-level intelligence, token-bounded chunks, and vector embeddings.
Extraction, normalization, chunking, and embedding are required for retrieval;
technical analysis and Document Understanding are separately observable
enrichments. This chapter follows the data in that order.

\section{Upload and Validation}
\label{sec:upload-validation}

An upload is accepted only in the context of an authenticated project. The
backend first derives the current user identity and verifies that the target
project is owned by that user. A nonexistent project and a project belonging
to another account are deliberately exposed through the same safe not-found
response. Consequently, an unauthorized caller cannot use document endpoints
to confirm whether another user's project or document exists. The ownership
test occurs before file validation and storage, so rejected callers cannot
cause files to be written outside their authorized workspace.

The accepted input surface is intentionally narrow. Table~\ref{tab:upload-formats}
summarizes the two supported formats and the checks applied by the server.
The size limit is exactly 20~MiB, or $20 \times 1024^2$ bytes. The multipart
request limit includes an additional 64~KiB allowance for form framing, but
this allowance does not increase the permitted file size.

\begin{table}[!htbp]
	\centering
	\small
	\caption{Accepted upload formats and validation checks}
	\label{tab:upload-formats}
	\begin{tabular}{p{2.0cm}p{3.0cm}p{4.4cm}p{4.0cm}}
		\toprule
		Format & Extension & Expected media type & Content validation \\
		\midrule
		PDF & \texttt{.pdf} & \texttt{application/pdf} & The first bytes must contain the PDF signature \texttt{\%PDF-}. \\
		DOCX & \texttt{.docx} & Open XML word-processing media type & The file must have a ZIP signature and contain both \texttt{[Content\_Types].xml} and \texttt{word/document.xml}. \\
		\bottomrule
	\end{tabular}
\end{table}

Validation is not based on the extension alone. The basename is trimmed,
limited to 255 characters, and rejected if it contains control characters.
Extension, media type, and the format checks in Table~\ref{tab:upload-formats}
must agree. The format-specific parser remains responsible for detecting
deeper damage, password protection, or unsupported content.

After validation, the stream is written to a local \texttt{Uploads} directory
beneath the server content root. A random identifier and the validated
lower-case extension form the storage name; the resolved path is checked to
remain inside that directory. The original name is stored separately for
display and citations. If database persistence fails, the controller attempts
to remove the newly written file.

The database record is initially marked \emph{Uploaded}. Related diagnostic
records are created at the same time: Document Understanding is \emph{Pending};
technical analysis and OCR are \emph{NotAnalyzed} for PDFs and \emph{Skipped}
for DOCX files. The identifier is then offered to a bounded in-memory channel.
The channel has capacity for 100 unique document identifiers, permits multiple
writers, and is consumed by one hosted background worker. Suppressing duplicate
identifiers prevents repeated button presses from scheduling the same document
multiple times. The worker creates a fresh dependency-injection scope for each
job and invokes the processing service outside the upload request.

The queue keeps the request responsive, but it is process-local and
non-durable. Pending items can be lost when the server stops, while persisted
\emph{Uploaded} or \emph{Failed} records can be queued again. A durable broker
and multi-worker coordination are future work.

Figure~\ref{fig:ingestion-pipeline} reserves the complete flow diagram. It
shows that validation and storage lead into one connected processing path,
rather than a set of unrelated document utilities.

\thesisfigureplaceholder{DIA-03}{Complete document-ingestion pipeline from validated upload to searchable persisted knowledge.}{fig:ingestion-pipeline}

\section{PDF and DOCX Text Extraction}
\label{sec:text-extraction}

The processing worker selects an extractor from the validated media type and
stored extension. PDF extraction is implemented with PdfPig, while DOCX
extraction uses the Open XML SDK. \citationneeded{PdfPig project and PDF text extraction}
\citationneeded{Open XML SDK and WordprocessingML specification}
The result of either adapter is the same application-level structure: an
ordered collection of text sections carrying optional page and heading
provenance plus an extraction-method marker.

\begin{table}[!htbp]
	\centering
	\small
	\caption{Format-specific extraction characteristics}
	\label{tab:extraction-comparison}
	\begin{tabular}{p{2.2cm}p{4.0cm}p{4.0cm}p{3.5cm}}
		\toprule
		Aspect & PDF & DOCX & Persisted provenance \\
		\midrule
		Primary order & Physical page order & Body-child order & Sequential section index \\
		Text unit & One non-empty native page & Heading-delimited group of paragraphs and tables & Raw text and extraction method \\
		Location & One-based page number & No reconstructed page number & Page when available \\
		Structure & Reading-order text supplied by PdfPig & Heading styles 1--3 and table row/cell separators & Section title when available \\
		\bottomrule
	\end{tabular}
\end{table}

For a PDF, PdfPig opens the source stream and visits pages in their original
order. Its content-order extractor reconstructs a text string from PDF text
objects. Leading and trailing whitespace is removed, and a section is emitted
only when the page contains non-empty native text. The section records the
one-based page number and the \emph{NativePdf} extraction method. Skipping an
empty native page is intentional: the subsequent OCR stage can add recognized
text when technical analysis identifies the page as a scan. Parser errors are
translated into a safe message explaining that the PDF may be damaged or
password protected, without returning low-level parser details to the user.

Because PDF text may consist of positioned glyphs rather than semantic
paragraphs, unusual layouts, columns, or font maps can still produce imperfect
reading order. Retained page provenance later allows chunks and answer sources
to identify where evidence originated.

For DOCX, the extractor opens the WordprocessingML package read-only and walks
the direct children of the main document body. Paragraph text is assembled
from its text descendants. A paragraph whose style identifier normalizes to
\emph{Heading1}, \emph{Heading2}, or \emph{Heading3} starts a new logical
section. The heading is kept in the content and is also stored as the section
title, truncated to the database limit where necessary. This dual treatment
is useful: the words remain searchable, while the title can also be displayed
as provenance and used as a preferred chunk boundary.

Tables are converted without inserting a large source-code-like representation.
Paragraphs within a cell are joined with spaces, cells are separated by a
vertical bar, and rows are separated by line endings. This preserves row order
and makes relationships in modest tables readable to later retrieval stages.
Paragraphs and tables remain in body order. The extractor does not reconstruct
DOCX pagination or explicitly traverse headers, footers, comments, footnotes,
drawings, or heading levels outside the recognized three. It therefore makes
no claim of complete WordprocessingML coverage.

Both extractors preserve their output as the raw representation; normalization
does not overwrite it. A later full reprocessing may atomically replace both
raw and normalized forms, so ``raw'' means the unaltered authoritative
extraction, not a history of every run.

\section{Raw and Normalized Text}
\label{sec:text-normalization}

Repeated page edges, divided words, and inconsistent whitespace consume chunk
budgets and can become misleading retrieval signals. Conversely, aggressive
cleanup can delete clauses or identifiers. The implementation therefore
prefers retaining occasional boilerplate to silently removing valid evidence.
\citationneeded{text normalization and conservative de-hyphenation}

The first normalization pass is format-independent. Sections are ordered by
their source index. Windows and classic Macintosh line endings are converted to
a single newline convention. Horizontal whitespace on each line is collapsed
to one space, line edges are trimmed, and consecutive empty lines are reduced
to one paragraph break. This gives deterministic text without flattening all
paragraph structure.

PDF page-edge cleanup is activated only when at least three page-numbered
sections exist. For each page, the first and last 15 non-empty lines form the
header and footer candidate regions. Individual repeated lines are considered
only near the physical edge, within the first or last three eligible positions.
Multi-line candidates may extend farther into the 15-line region, but their
boundary must begin or end no more than two lines away from the corresponding
page edge. This restriction reduces the risk of removing repeated phrases from
the body.

Candidate comparison is deliberately tolerant of harmless representation
differences. Text is normalized to Unicode compatibility form, whitespace is
collapsed, spaces around common punctuation are standardized, and comparison
is case-insensitive. Multi-line candidates are compared as logical blocks, so
different physical wrapping does not necessarily prevent recognition. A block
must contain at least 40 and at most 4,000 canonical characters. Normally it is
removed only when it occurs on at least the ceiling of 60\% of the eligible
pages. A long block of at least 160 characters can also be confirmed in a dense
local sequence: at least three occurrences must occupy a page span whose
density reaches the same 60\% threshold. The local rule handles, for example,
a repeated legal footer that appears in one substantial section of a long
combined document without lowering the global threshold.

Standalone page numbers receive separate treatment. Recognized forms include
a plain integer and forms using words such as ``page'', ``pagina'', ``of'', or
``din''. When the physical PDF page is known, the displayed current number must
equal it; a different number in the body is preserved. Patterns resembling
numbered headings, such as ``4.2 Processing architecture'', are protected from
repetition removal. These checks do not guarantee detection of every decorative
pagination scheme, but they prevent a general numeric-removal rule from
destroying meaningful quantities.

Word-break repair is similarly bounded. A line must end in one hyphen rather
than a double hyphen, the next line must begin with a lower-case letter, the
alphabetic prefix before the hyphen must contain at least four letters, and the
following alphabetic run must contain at least three. The next run must not
immediately continue into another hyphenated form. Only then is the line-ending
hyphen removed and the two parts joined. The rule repairs common PDF line-wrap
hyphenation while preserving many compound words and numeric identifiers. It is
language-independent at the character-class level and therefore works with
Romanian diacritics and other Unicode letters, although no dictionary is used.

\begin{table}[!htbp]
	\centering
	\small
	\caption{Normalization operations, purpose, and conservative boundary}
	\label{tab:normalization-operations}
	\begin{tabular}{p{3.2cm}p{4.4cm}p{5.4cm}}
		\toprule
		Operation & Intended benefit & Conservative boundary \\
		\midrule
		Line-ending and whitespace normalization & Stable text and token counts & Paragraph separation is retained. \\
		Repeated edge-line removal & Reduce common page headers and footers & Requires PDF pages, edge position, and 60\% repetition. \\
		Logical block removal & Handle multi-line boilerplate with changed wrapping & Length, edge-offset, global, and dense-local thresholds apply. \\
		Page-number removal & Remove isolated pagination noise & Restricted to page-edge regions and expected page number. \\
		Word-break repair & Rejoin words divided by line wrapping & Requires a specific alphabetic lower-case pattern. \\
		Whole-document fallback & Prevent destructive cleanup & Restores whitespace-normalized input if all meaningful text would vanish. \\
		\bottomrule
	\end{tabular}
\end{table}

After edge cleanup and optional word repair, whitespace is normalized once
more. The service records whether each section changed, its non-negative net
character reduction, and the normalization time. Aggregate counts on the
document support diagnostics. Most importantly, a whole-document guard checks
whether every normalized section became empty even though the original
contained meaningful non-page-number text. In that situation, the system
falls back to whitespace-normalized original sections. The result may contain
some unwanted repetition, but it does not lose the document's entire useful
content because of a heuristic.

The explicit rebuild reads stored raw sections and prepares normalization,
chunks, and embeddings before replacement. Failure preserves the previous
usable representation and records the appropriate error, making heuristic
cleanup observable and recoverable.

\section{Technical PDF Analysis}
\label{sec:technical-pdf-analysis}

PDFs may contain native text, scanned images, or a mixture. The application
therefore performs deterministic structural analysis before OCR routing. The
analysis describes page representation, not subject matter.
\citationneeded{PDF text and raster-image representation}

The controlled types are listed in Table~\ref{tab:technical-types}. Keeping the
taxonomy small makes routing and diagnostics understandable. In particular,
\emph{ImageBased} is not treated as a synonym for \emph{Scanned}: a page may
contain images without a single page-sized raster that justifies scan OCR.

\begin{table}[!htbp]
	\centering
	\small
	\caption{Technical PDF page types}
	\label{tab:technical-types}
	\begin{tabular}{p{2.6cm}p{10.4cm}}
		\toprule
		Type & Interpretation \\
		\midrule
		TextBased & A meaningful native text layer is the primary useful representation. \\
		Scanned & Little or no meaningful text exists and one raster placement covers at least 80\% of the visible page. \\
		ImageBased & Little or no meaningful text exists and raster content is present, but no accepted image reaches the page-sized threshold. \\
		Mixed & Meaningful text coexists with substantial raster coverage, or materially different known page types coexist at document level. \\
		Unknown & Structural evidence is insufficient for a safe classification. \\
		\bottomrule
	\end{tabular}
\end{table}

PdfPig supplies two groups of measurements for each page. The text scan counts
Unicode letters and digits and also counts alphanumeric word runs containing at
least two characters. Text is meaningful when either at least 40 alphanumeric
characters or at least eight useful word runs are present. The logical
disjunction is important: short tabular pages can have several meaningful
words without reaching 40 characters, while a page containing a long compact
identifier sequence can reach the character criterion.

Image masks and invalid sample dimensions are ignored. Each remaining bounding
box is intersected with the page crop box, and coverage is the largest visible
single-image area divided by page area, clamped to $[0,1]$. Image areas are not
summed, avoiding inflation from overlapping placements.

\begin{table}[!htbp]
	\centering
	\small
	\caption{Deterministic PDF-classification thresholds}
	\label{tab:pdf-thresholds}
	\begin{tabular}{p{7.8cm}r}
		\toprule
		Threshold & Configured value \\
		\midrule
		Minimum alphanumeric characters for meaningful text & 40 \\
		Minimum useful word runs for meaningful text & 8 \\
		Minimum alphanumeric length of a useful word run & 2 \\
		Substantial raster coverage & 0.30 \\
		Page-sized raster coverage & 0.80 \\
		Strong document-level majority & 0.80 \\
		Maximum tolerated unknown-page ratio & 0.20 \\
		Minimum unknown-page tolerance & 1 page \\
		\bottomrule
	\end{tabular}
\end{table}

Algorithm~\ref{alg:pdf-classification} expresses the actual decision order.
Notice that meaningful text takes precedence over the scanned-page branch. A
page with a page-sized image and a meaningful text layer is \emph{Mixed}, not
\emph{Scanned}; it is therefore not sent to OCR by the current routing policy.

\begin{algorithm}[!htbp]
	\caption{Deterministic technical PDF classification}
	\label{alg:pdf-classification}
	\begin{algorithmic}[1]
		\Require Ordered PDF pages with text and accepted raster placements
		\Ensure A technical type for every page and for the document
		\ForAll{pages in source order}
			\State Count alphanumeric characters and word runs of length at least 2
			\State $meaningful \gets (characters \geq 40) \lor (words \geq 8)$
			\State $coverage \gets$ largest visible single-image coverage, clamped to $[0,1]$
			\State $pageSized \gets (imageCount > 0) \land (coverage \geq 0.80)$
			\If{$meaningful$}
				\State $type \gets$ Mixed if $coverage \geq 0.30$, otherwise TextBased
			\ElsIf{$pageSized$}
				\State $type \gets$ Scanned
			\ElsIf{$imageCount > 0$}
				\State $type \gets$ ImageBased
			\Else
				\State $type \gets$ Unknown
			\EndIf
		\EndFor
		\State $unknownTolerance \gets \max(1,\lfloor 0.20 \times pageCount \rfloor)$
		\If{there are no known pages or unknown pages exceed the tolerance}
			\State \Return Unknown
		\EndIf
		\State Group known pages by type and find the stable dominant group
		\If{the dominant group represents at least 80\% of known pages}
			\State \Return the dominant type
		\EndIf
		\State \Return Mixed when several known types remain; otherwise Unknown
	\end{algorithmic}
\end{algorithm}

At document level, empty or all-unknown input is \emph{Unknown}. Unknown pages
may not exceed the greater of one and the floor of 20\% of all pages. Among
known pages, an inclusive 80\% majority wins; otherwise coexisting known types
produce \emph{Mixed}. Stable ordering makes ties deterministic.

The result stores page metrics and aggregate counts. Reuse requires the
original-file SHA--256 hash and analyzer version
\texttt{pdf-technical-analysis-v1} to match; a forced rebuild bypasses reuse.
Failure is recorded separately and native extraction continues.

The intended flow of these decisions is shown by the reserved
Figure~\ref{fig:pdf-classification-flow}.

\thesisfigureplaceholder{DIA-04}{Decision flow for deterministic page-level and document-level technical PDF classification.}{fig:pdf-classification-flow}

\section{Selective Local OCR}
\label{sec:selective-ocr}

OCR is applied selectively, after technical analysis and native extraction.
The routing policy chooses only pages whose technical type is exactly
\emph{Scanned}. It does not currently OCR \emph{ImageBased}, \emph{Mixed},
\emph{TextBased}, or \emph{Unknown} pages. This narrow rule avoids replacing a
usable text layer with a less certain recognition result and prevents every
decorative image from triggering expensive rasterization. It also makes the
routing decision reproducible from persisted page diagnostics.

Rendering is provided by PDFtoImage using its PDFium runtime assets, and
recognition is performed by Tesseract through the TesseractOCR .NET wrapper.
\citationneeded{PDFium and PDFtoImage project documentation}
\citationneeded{Tesseract OCR official project documentation}
Both operations execute on the application host; the OCR stage itself does not
send page images to a cloud OCR service. This statement does not imply that all
document data remains local. Later Document Understanding and embedding stages
may transmit bounded text to the configured model provider.

\begin{table}[!htbp]
	\centering
	\small
	\caption{Selective OCR configuration}
	\label{tab:ocr-configuration}
	\begin{tabular}{p{6.8cm}p{6.0cm}}
		\toprule
		Setting & Current repository value \\
		\midrule
		OCR enabled & Yes \\
		Tesseract language set & \texttt{ron+eng} \\
		Requested render resolution & 300 DPI \\
		Maximum scanned candidate pages & 200 \\
		Maximum rendered pixels per page & 25,000,000 \\
		Language-data directory & \texttt{tessdata} beneath the server content root \\
		Routing policy/version & Scanned pages only, \texttt{ocr-routing-v1} \\
		\bottomrule
	\end{tabular}
\end{table}

Before recognition, the service validates languages, limits, and the data path.
A configuration hash covers trained-data fingerprints, engine identity,
languages, DPI, and limits; a routing hash covers selected page-number/type
pairs. With the source hash, these identities control safe reuse.

For each selected page, pixel dimensions are calculated at 300 DPI. If they
exceed 25 million pixels, DPI is scaled by the square root of the permitted area
ratio and reduced further after rounding if needed. A white-background PNG is
rendered in memory. This bounds CPU and memory use while retaining the highest
resolution allowed by the rule.

Tesseract processes the image with the Romanian and English models. Line
endings are normalized, confidence is converted to the interval from zero to
one, and leading/trailing whitespace is removed. Non-empty output becomes an
extracted section with the original page number and \emph{Ocr} extraction
method. Empty output is recorded as an \emph{Empty} page result and does not
create artificial content. Render or recognition errors likewise create safe
diagnostics without inserting text. If an error indicates that the OCR
infrastructure itself is unavailable, remaining selected pages are marked
failed without repeatedly attempting the same unavailable resource.

Algorithm~\ref{alg:selective-ocr} summarizes routing, reuse, bounded execution,
and unification with native extraction.

\begin{algorithm}[!htbp]
	\caption{Selective local OCR routing and extraction}
	\label{alg:selective-ocr}
	\begin{algorithmic}[1]
		\Require PDF stream, native page sections, technical-analysis diagnostics, force flag
		\Ensure Ordered unified sections and persisted OCR diagnostics
		\If{technical analysis is not Ready}
			\State Persist OCR as NotAnalyzed and \Return reindexed native sections
		\EndIf
		\State $candidates \gets$ Scanned pages ordered by page number
		\If{$candidates$ is empty}
			\State Persist OCR as Skipped and \Return native sections
		\EndIf
		\State Validate configuration and local language data
		\If{configuration or infrastructure is unavailable}
			\State Persist safe Failed diagnostics
			\State \Return native sections excluding current scanned candidate pages
		\EndIf
		\State Compute source, routing, and configuration identities
		\If{not forced and a complete matching Ready result plus stored OCR text exists}
			\State \Return native and reused OCR sections in page order
		\EndIf
		\State Select at most the first 200 candidate pages
		\ForAll{selected candidate pages}
			\State Render at the highest DPI that satisfies the 25,000,000-pixel limit
			\State Run local Tesseract with \texttt{ron+eng}
			\State Persist Ready text, Empty, or safe Failed page diagnostics
		\EndFor
		\State Derive Ready, Partial, or Failed aggregate status
		\State Remove native sections for all scanned candidate pages
		\State Add successful OCR sections, sort by page, and assign new section indexes
		\State \Return the unified sections
	\end{algorithmic}
\end{algorithm}

The replacement step deserves clarification. Native sections from pages
classified as scanned are excluded and only successful OCR output for those
pages is inserted. This avoids combining a few noisy native glyphs with the
recognized page, but it also means that a failed scanned page contributes no
text. Native sections from non-candidate pages are kept. Thus a mixed PDF can
still become \emph{Ready} when one scan fails but other native or recognized
pages provide useful content; an all-scanned document for which every page
fails OCR has no extractable text and the overall ingestion correctly fails.

Zero candidates yields \emph{Skipped}; all successes yield \emph{Ready}; some
successes yield \emph{Partial}; and none yields \emph{Failed}. Page states are
\emph{Ready}, \emph{Empty}, \emph{Failed}, and \emph{SkippedLimit}. Only the
first 200 candidates are processed, but all count in the aggregate. The first
excess page receives \emph{SkippedLimit}; further excess pages do not each
create a diagnostic row.

Stored provenance includes all three hashes, engine/languages, limits,
effective DPI and dimensions, duration, recognized counts, confidence, source
type, and whether text entered extraction. Reuse requires a completely
\emph{Ready} matching result and one stored OCR section per candidate. Partial
or failed results are rerun, and force always bypasses reuse.

The TesseractOCR wrapper dependency is version 5.5.2, while the service has a
5.5.1 fallback engine label and can later record the native engine's value.
Because package and native-engine versions differ conceptually, the thesis says
Tesseract~5 through the wrapper until the submission runtime is verified.

Figure~\ref{fig:selective-ocr-flow} will visualize the page-level route and the
return to the unified extraction sequence.

\thesisfigureplaceholder{DIA-05}{Selective OCR flow from Scanned-page diagnostics through bounded PDFium rendering and local Tesseract recognition.}{fig:selective-ocr-flow}

\section{Document Understanding}
\label{sec:document-understanding}

Once normalized text exists, the application attempts a bounded
document-level analysis. This stage does not answer a user's question and does
not replace retrieval evidence. Its purpose is to produce a concise structured
description that can support document diagnostics and later act as a soft
ranking signal. The current configured model is \texttt{gpt-5.6-luna}; if a
separate Understanding model is not configured, the service falls back to the
answer model. Provider interaction uses structured output rather than parsing
free-form prose. \citationneeded{OpenAI Responses API structured outputs}

The output contains a controlled document type, optional subtype, confidence,
primary language code and confidence, optional detected title and subject, and
an ordered collection of generic metadata entries. The document-type allowlist
is: Unknown, Contract, Invoice, Receipt, Report, Policy, Procedure, Manual,
CourseMaterial, ResearchPaper, FinancialDocument, Form, Letter, Resume,
TechnicalDocument, and Other. This taxonomy is intentionally broad enough for
the project corpus without presenting an open-ended model-generated category
as an authoritative database value.

\begin{table}[!htbp]
	\centering
	\small
	\caption{Controlled Document Understanding metadata kinds}
	\label{tab:understanding-metadata}
	\begin{tabular}{p{3.0cm}p{9.8cm}}
		\toprule
		Kind & Intended information \\
		\midrule
		Organization & Companies, institutions, public bodies, and other named organizations \\
		Person & People explicitly named by the document \\
		Identifier & Contract numbers, invoice identifiers, references, registration codes, and similar values \\
		Date & Explicit relevant dates, conservatively normalized when unambiguous \\
		MonetaryAmount & Amounts and currencies as supported by document text \\
		Jurisdiction & Explicit territorial or legal jurisdiction \\
		Topic & Concise subject keywords supported by the text \\
		Other & Useful generic metadata that does not fit the preceding kinds \\
		\bottomrule
	\end{tabular}
\end{table}

The model receives explicit higher-priority instructions that document
contents are untrusted data, never instructions. It is told not to follow
requests embedded in the document, reveal prompts or secrets, call tools,
contact people, or invent facts. The user message surrounds document content
with visible begin/end delimiters. The JSON schema prohibits additional fields,
enumerates every accepted document and metadata kind, requires all principal
fields, bounds confidence values, and limits metadata to 50 entries. Provider
output storage is disabled for this request. These measures reduce the attack
surface of indirect prompt injection, but they do not justify a claim that the
general problem has been completely solved.

Structured output is followed by local validation. Enum text must map to an
allowlisted value. Required confidence scores and optional metadata confidence
must be finite and between zero and one; invalid values fail the whole attempt
rather than being silently clamped. Language tags are checked for a
BCP-47-compatible shape and deterministically cased. Subtype, title, subject,
label, and value lengths are bounded. Labels are collapsed into ASCII
lower-snake-case form. Duplicate metadata entries are removed deterministically
using kind, label, and normalized or original value. Dates are converted to
\texttt{yyyy-mm-dd} only when a supported representation is unambiguous,
including selected Romanian and English month-name forms. Ambiguous numeric
dates remain without a normalized value. Monetary values are preserved as
cleaned source strings; the local code does not invent a numeric/currency
normalization.

Sending every token of a very long document would be costly and would violate
the bounded-input requirement. The input builder therefore orders non-empty
normalized sections, joins their canonical content with paragraph breaks, and
computes a SHA--256 hash over this entire canonical text. A separate annotated
form adds page and heading markers for the model. An input is usable only when
it contains at least 20 tokens and at least 40 letters or digits. Shorter input
is deterministically marked \emph{Skipped} and no provider call is made.

If the annotated form contains no more than 6,000 \texttt{cl100k\_base} tokens,
it is used in full. Otherwise a deterministic representative sample is built.
After accounting for the three sample markers and a 128-token safety reserve,
one half of the available budget is assigned to the beginning, one quarter to
a centered middle window, and the remaining quarter to the end. The middle
window takes a suffix from the left side and a prefix from the right side of the
central region. Token boundaries are clamped so a UTF-16 surrogate pair is not
split. The sample therefore represents opening context, central body material,
and conclusions while remaining reproducible. It is representative rather
than exhaustive; facts located only outside these windows may not become
metadata, even though they remain available to chunk retrieval.

Idempotency uses the full canonical-content hash, configured model, and prompt
version \texttt{document-understanding-v1}. A matching \emph{Ready} or
\emph{Skipped} result is reused unless the caller explicitly forces analysis.
Because the hash covers the full text rather than the sample, a change in an
unsampled region still invalidates the previous result. Successful metadata
replacement is transactional. Model or validation failure marks the attempt
\emph{Failed} with a safe message and does not persist a partial provider
response.

Most importantly, Understanding failure is non-fatal to ordinary RAG
ingestion. The processing service logs the safe failure and continues with the
same normalized sections, chunking, and embedding. After normalization or OCR
rebuild, a changed input is first staged as \emph{Pending}; the new searchable
representation is committed, and Understanding is retried afterward. If that
retry fails, the document remains \emph{Ready} for retrieval. This ordering
reflects the role of Understanding as useful enrichment rather than the sole
source of searchable content.

\section{Token-Aware Document Chunking}
\label{sec:document-chunking}

The normalized document is usually too large and too heterogeneous to serve as
one retrieval item. A single vector for an entire report would blend unrelated
topics, while inserting the full document into every answer request would
exceed bounded context and make provenance imprecise. At the opposite extreme,
one vector per sentence can remove the surrounding definitions and conditions
needed to interpret a passage. Chunking establishes an intermediate evidence
unit: small enough to rank and cite, but large enough to preserve local meaning.
Overlap carries limited context across boundaries. \citationneeded{document chunking for semantic retrieval}

The implementation measures tokens with the \texttt{cl100k\_base} encoding
provided through Microsoft.ML.Tokenizers. \citationneeded{cl100k base tokenizer and Microsoft ML Tokenizers}
Character counts are retained for diagnostics, but all size decisions are made
from actual tokenizer counts. This is important for multilingual material,
punctuation, numbers, and Unicode because equal character lengths do not imply
equal model-token lengths.

\begin{table}[!htbp]
	\centering
	\small
	\caption{Current chunking configuration and semantics}
	\label{tab:chunking-configuration}
	\begin{tabular}{p{3.4cm}r p{7.2cm}}
		\toprule
		Parameter & Tokens & Meaning \\
		\midrule
		Target & 700 & Preferred point for closing a base plan when a natural unit would take it over the target. \\
		Hard maximum & 900 & Absolute limit checked against every emitted chunk, including overlap. \\
		Overlap & 100 & Maximum suffix budget copied from the preceding base plan when space permits. \\
		Minimum useful size & 100 & Preferred lower bound before heading/target closure and for final-tail rebalancing. \\
		\bottomrule
	\end{tabular}
\end{table}

The target and minimum require careful interpretation. They are planning goals,
not unconditional assertions about every possible output. A natural text unit
that contains 800 tokens is retained because it is below the 900-token maximum,
even though it exceeds the 700-token target. A very short document remains one
short chunk. Likewise, the final chunk may remain below 100 tokens when moving
or merging available units would violate another constraint. The 900-token
maximum is the hard invariant.

\subsection{Creating provenance-carrying text units}

Source sections are sorted by their persisted index. Each section is split on
line endings; empty entries are discarded and non-empty entries are trimmed.
Every resulting unit carries its source-section index, optional page number,
optional section title, and a flag indicating whether it is the first unit of a
titled section. Separators are also explicit: ordinary source lines are joined
as paragraphs, while sentences cut from one oversized line are rejoined with a
space. These details mean that token counts are computed over the same text that
will actually be persisted, including separators.

A unit at or below 900 tokens is left intact. Only an oversized unit is divided
into sentences. The sentence detector considers a full stop, exclamation mark,
or question mark followed by whitespace or common closing punctuation. It does
not split a decimal point between digits, and it protects a compact allowlist of
common Romanian and English abbreviations such as \emph{art.}, \emph{nr.},
\emph{dl.}, \emph{dna.}, \emph{dr.}, \emph{prof.}, \emph{etc.}, \emph{Mr.},
\emph{e.g.}, and \emph{i.e.}. This is a conservative boundary heuristic, not a
general linguistic sentence parser.

If one sentence still exceeds 900 tokens, the tokenizer locates the largest
prefix within the limit. The splitter moves backward to whitespace when doing
so produces a non-empty prefix; otherwise it uses the exact token boundary. It
also checks UTF-16 boundaries and moves away from the middle of a surrogate
pair. Whitespace at the start of the remaining suffix is consumed, and the
process repeats. Thus even an unusually long uninterrupted passage can be
bounded without corrupting supplementary Unicode characters.

\subsection{Building and rebalancing base plans}

Base plans are constructed by scanning units once in source order. A new titled
section closes the current plan only after that plan contains at least 100
tokens. This preference avoids attaching a substantial previous section to a
new heading while also preventing every tiny heading fragment from creating an
undersized chunk. Before adding each unit, the generator calculates the exact
candidate token count. If the candidate would exceed the 700-token target and
the current plan is already at least 100 tokens, the current plan closes. If a
candidate would exceed 900 tokens under any remaining circumstance, the current
non-empty plan closes regardless of its size. Since every individual unit has
already been bounded, the newly started plan remains within the hard maximum.

The last plan receives one additional balancing pass. If it already contains
at least 100 tokens, nothing changes. Otherwise, the generator first attempts
to merge it with the preceding plan; this is accepted only if their exact joined
content remains within 900 tokens. If they cannot be merged, complete units are
moved from the end of the preceding plan to the beginning of the final plan.
Each move must leave the preceding plan at or above 100 tokens and the final
plan at or below 900. Moving whole units preserves natural boundaries. The pass
stops when the tail reaches the desired minimum or no valid move exists.

\subsection{Adding bounded overlap}

Overlap is calculated after base planning, so it cannot recursively include
content that was itself copied as overlap. For every plan after the first, the
available overlap budget is the smaller of 100 tokens and the space remaining
under 900 after counting the current base plan. The algorithm walks backward
through the immediately preceding base plan and prepends complete trailing
units while they fit. If the first considered unit is larger than the entire
budget and no complete unit has been selected, it takes a tokenizer-bounded
suffix of that unit and advances to a safe word start.

The generator explicitly prevents the overlap from reproducing the entire
previous base plan. If several selected units would do so, the earliest is
removed. If the previous plan consists of one unit, only a reduced suffix is
attempted. After overlap and base units are combined, leading overlap units are
discarded if exact counting still exceeds 900. The base content itself is never
removed to make room for overlap. As a result, overlap is beneficial but
subordinate to the hard size bound and to preservation of the current evidence.

Algorithm~\ref{alg:token-aware-chunking} presents these operations in a compact
form while preserving the decisions made by the production implementation.

\begin{algorithm}[!htbp]
	\caption{Token-aware document chunking with bounded overlap}
	\label{alg:token-aware-chunking}
	\begin{algorithmic}[1]
		\Require Ordered normalized sections with page and heading provenance
		\Ensure Ordered non-empty chunks of at most 900 tokens
		\State $units \gets$ empty list
		\ForAll{sections ordered by source index}
			\ForAll{non-empty trimmed lines in the section}
				\If{the line contains at most 900 tokens}
					\State Append one provenance-carrying unit
				\Else
					\State Split at protected sentence boundaries
					\State Split any oversized sentence at a safe token/word/Unicode boundary
				\EndIf
			\EndFor
		\EndFor
		\If{$units$ is empty}
			\State Fail with a safe no-text error
		\EndIf
		\State $plans \gets$ empty list; $current \gets$ empty plan
		\ForAll{$unit$ in source order}
			\If{$unit$ starts a titled section and $current \geq 100$ tokens}
				\State Close $current$ and begin an empty plan
			\EndIf
			\If{$current + unit > 700$ tokens and $current \geq 100$ tokens}
				\State Close $current$ and begin an empty plan
			\EndIf
			\If{$current$ is non-empty and $current + unit > 900$ tokens}
				\State Close $current$ and begin an empty plan
			\EndIf
			\State Append $unit$ to $current$
		\EndFor
		\State Close the final non-empty plan
		\If{the final plan is below 100 tokens and a previous plan exists}
			\State Merge when the result is at most 900 tokens
			\State Otherwise move valid whole trailing units from the previous plan
		\EndIf
		\For{$i \gets 1$ to number of plans}
			\State $base \gets plans[i]$
			\If{$i > 1$}
				\State $budget \gets \min(100, 900 - \Call{TokenCount}{base})$
				\State $overlap \gets$ largest safe, non-whole suffix of $plans[i-1]$ within $budget$
			\Else
				\State $overlap \gets$ empty
			\EndIf
			\State Remove leading overlap units while $overlap + base > 900$ tokens
			\State Build exact content and count its tokens
			\State Reject empty or over-limit output
			\State Emit index, content, page range, title, and source-section range
		\EndFor
	\end{algorithmic}
\end{algorithm}

\subsection{Persisted provenance and trade-offs}

For an emitted chunk, page start and end are derived from all combined units,
including overlap. Source-section bounds likewise use the minimum and maximum
source indexes represented in the chunk. The displayed section title is the
first non-empty title in the current base plan, not a title copied only through
overlap. This prevents an inherited suffix from misnaming the principal new
content. Chunks may span source sections and PDF pages; page boundaries are not
hard stops. Their page range therefore describes all evidence included in the
text rather than pretending that every chunk belongs to one page.

\begin{table}[!htbp]
	\centering
	\small
	\caption{Chunk-size and overlap trade-offs}
	\label{tab:chunking-tradeoffs}
	\begin{tabular}{p{3.1cm}p{4.7cm}p{4.7cm}}
		\toprule
		Choice & Benefit & Cost or risk \\
		\midrule
		Very small chunks & Precise local matching and short context entries & Definitions and conditions can be separated; many vectors are required. \\
		Very large chunks & More surrounding context remains together & Embeddings mix topics, ranking becomes less precise, and context is consumed quickly. \\
		No overlap & Minimal duplication and storage & Relevant meaning can be split exactly at a boundary. \\
		Bounded overlap & Carries nearby context into the next evidence unit & Repeated text can appear in adjacent results and consumes part of the hard limit. \\
		Natural-unit preference & Preserves lines, sentences, and headings where possible & Output sizes vary around the preferred target. \\
		\bottomrule
	\end{tabular}
\end{table}

The algorithm is deterministic for identical normalized sections, metadata,
tokenizer, and options. It does not depend on random choices or a model call.
Repeated rebuilds therefore produce the same ordered chunk contents and
provenance, although database identifiers and timestamps can change when rows
are replaced. Romanian and English text, diacritics, and supplementary Unicode
characters are preserved. Determinism supports testing, reproducible retrieval
experiments, and reliable invalidation of downstream embeddings.

Figure~\ref{fig:chunk-overlap} will show two adjacent chunks, their shared
suffix, and the page/source metadata that accompanies both representations.

\thesisfigureplaceholder{DIA-06}{Conceptual token-bounded chunks with non-recursive overlap and preserved page, heading, and source-section provenance.}{fig:chunk-overlap}

\section{Embedding Generation}
\label{sec:embedding-generation}

Each completed chunk is mapped to one numerical embedding in the same order in
which the chunks were generated. The current provider model is
\texttt{text-embedding-3-small}, configured to return 1,536 dimensions.
Inputs are sent in sequential batches of at most 32. The adapter requests the
configured dimension explicitly, restores provider results to input order by
their response indexes, and rejects missing, duplicated, non-finite, or
wrong-sized vectors. \citationneeded{OpenAI text embedding model and dimensions}

The application persists vectors through the pgvector .NET type. PostgreSQL
enables the \texttt{vector} extension and stores the value in a native
\texttt{vector(1536)} column. \citationneeded{pgvector official documentation}
The fixed dimension is part of the database schema, so startup validation
rejects a different configured dimension rather than allowing a later database
error. Alongside the vector, every chunk stores the model identifier,
dimensions, generation time, and a SHA--256 hash of its exact content. The
document stores matching aggregate model, dimension, timestamp, and embedded
chunk count. These fields allow the API to distinguish a complete current
embedding set from stale or partially populated data without exposing raw
vectors to the frontend.

Currency validation is deliberately strict. The document must be \emph{Ready};
its embedded count must equal its chunk count; model and dimensions must match
current configuration; and every chunk must contain a vector, share the
document completion timestamp, and carry the correct hash for its present
content. A changed chunk therefore invalidates the aggregate even if older
document-level fields still look current.

There is no embedding cache that suppresses an explicit rebuild. The rebuild
operation intentionally regenerates vectors even when the existing set is
current. It first snapshots chunk identifiers, indexes, and exact content,
performs the provider call, reloads the rows, and verifies that the snapshot did
not change while generation was in progress. Only then are all vectors replaced
transactionally. If generation or validation fails, previous authoritative
embeddings remain available, the document returns to \emph{Ready}, and a safe
embedding error is recorded. Normalization and chunk rebuilds likewise prepare
new embeddings before replacing existing derived data.

The current schema does not define an HNSW or IVFFlat approximate-nearest-
neighbour index. It uses native pgvector representation and database-side
vector operations without claiming an approximate index that is not present.
This is suitable for the present local project scale and keeps the first
implementation straightforward. An approximate index, tuning process, and
larger-scale benchmark belong to future work.

Embedding is an external AI stage. Although OCR rendering and recognition are
local, chunk text is transmitted to the configured OpenAI endpoint to generate
vectors. API credentials are read from server configuration and are neither
stored with document records nor returned through document contracts.

\section{Status and Failure Isolation}
\label{sec:ingestion-status-isolation}

The pipeline separates overall usability from diagnostic enrichment. This
prevents one status value from hiding which stage failed and avoids making
optional analysis a prerequisite for all retrieval. Table~\ref{tab:processing-statuses}
summarizes the persisted concerns.

\begin{table}[!htbp]
	\centering
	\small
	\caption{Separation of ingestion and analysis status}
	\label{tab:processing-statuses}
	\begin{tabular}{p{3.0cm}p{5.2cm}p{4.8cm}}
		\toprule
		Concern & States or diagnostic fields & Role in ordinary ingestion \\
		\midrule
		Document processing & Uploaded, Processing, Ready, Failed & Overall searchable usability \\
		Technical PDF analysis & NotAnalyzed, Processing, Ready, Failed, Skipped & Non-fatal structural diagnostic and OCR prerequisite \\
		OCR document & NotAnalyzed, Processing, Ready, Partial, Failed, Skipped & Page recovery; non-fatal only when useful text remains \\
		OCR page & Ready, Empty, Failed, SkippedLimit & Per-page recognition and provenance \\
		Document Understanding & NotAnalyzed, Pending, Processing, Ready, Failed, Skipped & Non-fatal structured enrichment \\
		Normalization, chunks, embeddings & Counts, timestamps, and separate safe error fields & Required derived representation during initial ingestion \\
		\bottomrule
	\end{tabular}
\end{table}

A representative mixed-PDF outcome illustrates this design. Suppose one page
has native text, two pages are classified as scans, one scan is recognized, and
the other OCR attempt fails. OCR is \emph{Partial}, yet the unified extraction
contains the native page and one OCR page. If normalization, chunking, and
embedding succeed, the document can be \emph{Ready}. Diagnostics still reveal
the missing page, so readiness does not falsely imply perfect OCR coverage. If
all pages were scans and none produced text, the no-text check would instead
make document processing \emph{Failed}.

Similarly, a model-provider or validation problem can leave Document
Understanding \emph{Failed} while chunks and embeddings make the document
searchable. Technical analysis can fail while native PDF extraction succeeds;
in that case OCR cannot be safely routed and is reported as not analyzed, but
the native text can still proceed. These outcomes are not swallowed errors.
They are explicit, independently inspectable states.

Required-stage failures have stricter semantics during initial ingestion.
Extraction with no useful result, destructive-normalization prevention that
cannot produce valid input, chunk-generation failure, or embedding failure
prevents a new \emph{Ready} representation. Partial new text and chunks are
removed, aggregate counts are cleared, the stored source file is retained for
retry, and the error is assigned to the most specific field available. For
manual rebuilds of an already ready document, new derived data is prepared
before replacement; failure preserves the previous authoritative data and
records the rebuild error. This distinction lets the application be strict
when first establishing searchability and conservative when updating an
already usable document.

The final ingestion result is therefore more than extracted text. It is an
ordered and normalized representation, page and heading provenance,
independent technical and semantic diagnostics, deterministic bounded chunks,
validated vectors, and explicit status information. Chapter~\ref{chap:retrieval-grounding} uses these
persisted artifacts as the starting point for hybrid retrieval and grounded
question answering.
